% FORMALIZACION_NUEVA de la justificación del límite utilizado en
% 70_radiacion_termica.tex. No introduce un generador de constantes.
\subsection{Límite termodinámico de las sumas de ocupación}
\label{v:subsec:limite-termodinamico}

La densidad radial anterior puede obtenerse como límite efectivo de las
sumas de la realización periódica. La estimación debe conservar el modo
constante fijado y controlar uniformemente tanto el origen como la cola
espectral. Las representaciones de acción y calendario permanecen previas
a esta operación de realización.

\begin{lemma}[Convergencia de los tres observables térmicos]
\label{v:lem:limite-termodinamico}
Sean \(a,b>0\), \(L>0\), \(\Delta=2\pi/L\) y
\[
 F_N(r)=\frac{1}{e^{ar}-1},\qquad
 F_E(r)=\frac{br}{e^{ar}-1},\qquad
 F_Z(r)=-\log(1-e^{-ar})\quad(r>0).
\]
Para cada una de estas funciones se cumple
\begin{equation}
 \lim_{L\to\infty}\frac{2}{L^3}
 \sum_{\nu\in\mathbb Z^3\setminus\{0\}}
 F\!\left(\frac{2\pi}{L}|\nu|\right)
 =\frac{1}{\pi^2}\int_0^\infty r^2F(r)\,dr.
 \label{v:eq:limite-modos-discretos}
\end{equation}
Las sumas y la integral son finitas. La especialización
\(a=\beta\hbar c\), \(b=\hbar c\) proporciona, respectivamente,
las densidades de número, energía de excitación y logaritmo de la función
de partición de los modos transversales.
\end{lemma}

\begin{proof}
Las tres funciones son continuas y positivas en \((0,\infty)\).
Existen constantes \(C_0,C_1,d>0\), dependientes de \(a,b\) pero
independientes de \(\Delta\), tales que
\begin{equation}
 F(r)\leq\frac{C_0}{r}\quad(0<r\leq1),\qquad
 F(r)\leq C_1e^{-dr}\quad(r\geq1).
 \label{v:eq:envolventes-termicas}
\end{equation}
En efecto, \(e^{ar}-1\geq ar\) controla \(F_N\) y \(F_E\)
cerca del origen; además,
\(F_Z(r)=\log(1+F_N(r))\leq F_N(r)\leq1/(ar)\).
Para \(r\geq1\), la desigualdad
\(-\log(1-y)\leq y/(1-y)\), con \(y=e^{-ar}\), da la cota
exponencial de \(F_Z\). Las otras dos siguen de
\((e^{ar}-1)^{-1}\leq e^{-ar}/(1-e^{-a})\), absorbiendo el factor
\(r\) de \(F_E\) en una exponencial de menor tasa. Estas cotas
prueban ya la integrabilidad de \(r^2F(r)\).

Para las sumas discretas, agrupamos los puntos enteros por
\(m=\|\nu\|_\infty\geq1\). Cada capa contiene exactamente
\[
 (2m+1)^3-(2m-1)^3=24m^2+2
\]
puntos y satisface \(m\leq|\nu|\leq\sqrt3m\).
Fijemos \(0<\delta\leq1\) y \(0<\Delta\leq1\).
Si \(\delta<\Delta\), no hay puntos con
\(0<\Delta|\nu|\leq\delta\). En otro caso, poniendo
\(N=\lfloor\delta/\Delta\rfloor\), la primera cota de
\eqref{v:eq:envolventes-termicas} permite escribir
\begin{align*}
 \Delta^3\!\sum_{0<\Delta|\nu|\leq\delta}F(\Delta|\nu|)
 &\leq C_0\Delta^2\sum_{m=1}^{N}\frac{24m^2+2}{m}\\
 &\leq26C_0\Delta^2\sum_{m=1}^{N}m
 \leq26C_0\delta^2.
\end{align*}
En la segunda línea se ha usado \(m^{-1}\leq m\) y, después,
\(N(N+1)/2\leq N^2\) para \(N\geq1\).
La contribución integral del mismo entorno también es de orden
\(\delta^2\), pues
\(\int_0^\delta r^2F(r)\,dr\leq C_0\delta^2/2\).

Para \(R\geq1\), todo punto con \(\Delta|\nu|>R\) pertenece
a una capa \(m>R/(\sqrt3\Delta)\). La segunda cota de
\eqref{v:eq:envolventes-termicas} da
\begin{align*}
 \Delta^3\!\sum_{\Delta|\nu|>R}F(\Delta|\nu|)
 &\leq C_1e^{-dR/(2\sqrt3)}
 \Delta^3\sum_{m\geq1}(24m^2+2)e^{-d\Delta m/2}\\
 &\leq C_2e^{-dR/(2\sqrt3)},
\end{align*}
donde \(C_2\) no depende de \(0<\Delta\leq1\). Para comprobar
esa uniformidad, se toma \(q=e^{-d\Delta/2}\) y se utilizan
\[
 \sum_{m\geq1}q^m=\frac{q}{1-q},\qquad
 \sum_{m\geq1}m^2q^m=\frac{q(1+q)}{(1-q)^3}.
\]
La primera identidad procede de la serie geométrica y la segunda de
aplicar dos veces \(q\,d/dq\), operación válida en \(|q|<1\).
Además,
\(1-e^{-d\Delta/2}\geq\Delta(1-e^{-d/2})\) por concavidad.
El factor \(\Delta^3\) compensa, por tanto, los denominadores de
las dos expresiones. La cola integral tiende asimismo a cero por la
cota exponencial. Para cada \(\Delta>0\), estas mismas estimaciones
prueban la convergencia de la suma infinita.

En el anillo compacto \(\delta\leq|k|\leq R\), la función
\(F(|k|)\) es continua y acotada. Las sumas de Riemann de malla
\(\Delta\) convergen a su integral; las dos esferas de frontera
tienen medida nula. Se hace primero \(\Delta\to0\), después
\(\delta\to0\) y finalmente \(R\to\infty\), usando las cotas
uniformes anteriores. Resulta
\[
 \Delta^3\sum_{\nu\ne0}F(\Delta|\nu|)
 \longrightarrow \int_{\mathbb R^3}F(|k|)\,d^3k
 =4\pi\int_0^\infty r^2F(r)\,dr.
\]
Como \(2/L^3=2\Delta^3/(2\pi)^3\), el coeficiente radial final
es \(2\cdot4\pi/(2\pi)^3=1/\pi^2\), lo que demuestra
\eqref{v:eq:limite-modos-discretos}.
\end{proof}

El control del origen se efectúa sobre excitaciones de energía positiva.
El modo constante permanece fijado antes de formar la función de partición:
su eliminación posterior de una integral no sustituiría esa condición de
la representación finita. La prueba establece el paso al límite de la
realización declarada y conserva la distinción entre esta realización y
la construcción anterior de sus parámetros de acción y calendario.
